\documentclass[11pt,letterpaper]{article}
\usepackage[T1]{fontenc}
\usepackage[utf8]{inputenc}
\usepackage{lmodern}
\usepackage[margin=0.65in]{geometry}
\usepackage{xcolor}
\usepackage{hyperref}
\usepackage{fancyhdr}
\renewcommand{\familydefault}{\sfdefault}
\definecolor{ink}{HTML}{202A26}
\definecolor{muted}{HTML}{505A55}
\hypersetup{colorlinks=true,urlcolor=ink,pdftitle={Wanglong Lu - Applied Scientist Resume},pdfauthor={Wanglong Lu}}
\ifdefined\pdfgentounicode\input{glyphtounicode}\pdfgentounicode=1\fi
\setlength{\parindent}{0pt}
\setlength{\parskip}{3pt}
\setlength{\headheight}{14pt}
\pagestyle{fancy}
\fancyhf{}
\renewcommand{\headrulewidth}{0pt}
\fancyfoot[L]{\footnotesize\color{muted}Wanglong Lu}
\fancyfoot[R]{\footnotesize\color{muted}\thepage\ / 2}
\newcommand{\cvsection}[1]{\par\addvspace{9pt}{\large\bfseries #1}\par\vspace{1pt}\hrule\vspace{4pt}}
\newcommand{\role}[3]{\par\addvspace{6pt}\textbf{#1}\hfill{\small #2}\par{\small\color{muted}#3}\par}
\newenvironment{points}{\begin{list}{\textbullet}{\setlength{\leftmargin}{12pt}\setlength{\itemsep}{2pt}\setlength{\parsep}{0pt}\setlength{\topsep}{2pt}}}{\end{list}}
\newcommand{\project}[2]{\par\addvspace{5pt}\textbf{#1}\hfill{\small #2}\par}
\begin{document}
\fontsize{10.5}{12.7}\selectfont
\raggedright
{\fontsize{25}{28}\selectfont\bfseries Wanglong Lu}\par
{\large Senior Data Scientist \textbar\ Applied Multimodal Intelligence}\par
St. John's, NL, Canada \textbar\ \href{mailto:lwlxhl@gmail.com}{lwlxhl@gmail.com}\par
\href{https://longlongaaago.github.io/}{Website}\quad
\href{https://github.com/LonglongaaaGo}{GitHub}\quad
\href{https://www.linkedin.com/in/wanglong-lu-8a98b6235/}{LinkedIn}\quad
\href{https://scholar.google.com/citations?user=TuxCf4UAAAAJ}{Google Scholar}

\cvsection{Profile}
Applied ML researcher and Senior Data Scientist at Nasdaq, with a Ph.D. in Computer Science. Works on financial ML, multimodal image editing, and efficient adaptation, with experience deploying ML workflows on AWS. First-author \textbf{IEEE TNNLS and TVCG} publications; \textbf{NeurIPS} research integrated into \textbf{Hugging Face PEFT}.

\cvsection{Technical Skills}
\textbf{Core:} Python, PyTorch, SQL, Git, Bash, C++, CUDA; experimental design and feature engineering.\par
\textbf{ML \& GenAI:} Latent diffusion, Transformers, multimodal AI, PEFT/LoRA, Hugging Face, RAG, agentic tools.\par
\textbf{ML systems:} FastAPI, Mangum, PostgreSQL; AWS SageMaker, Bedrock, Lambda, S3, Docker, SLURM.

\cvsection{Professional Experience}
\role{Senior Data Scientist, AI/Analytics -- Nasdaq}{Mar. 2025 -- Present}{St. John's, NL, Canada}
\begin{points}
\item Designed cross-scale ensemble feature selection for \textbf{cheque-fraud ML (CFML)}, removing \textbf{77.2\% of input features} while retaining competitive performance in evaluations.
\item Deployed and integrated an \textbf{existing UI/backend prototype} for a cross-team ML training platform into an internal workspace; configured networking and IAM access, added \textbf{PostgreSQL-backed user-state persistence}, and deployed \textbf{FastAPI} services on \textbf{AWS Lambda using Mangum}.
\item Built a Python job-submission workflow that packages local scripts and dependencies for \textbf{SageMaker Processing, Training, and Hyperparameter Tuning}; linked jobs with \textbf{SageMaker Pipelines}.
\item Developed a \textbf{Bedrock-based feature-formula validator} with structured logical-consistency checks and reliability scores to flag hard-coded errors in feature-engineering formulas.
\item Built a \textbf{RAG-based incident-response assistant} as an internal hackathon project, retrieving troubleshooting context for on-call alerts through the company's GenAI platform.
\end{points}

\role{AI Algorithm Intern -- Nasdaq Verafin}{May 2024 -- Jan. 2025}{St. John's, NL, Canada}
\begin{points}
\item Designed and benchmarked diffusion-based restoration and inpainting for degraded cheque images using Python and PyTorch, comparing with existing restoration methods.
\item Developed patch-level Gaussian latent modeling for image verification; evaluated the approach and presented results to cross-functional colleagues.
\end{points}

\role{AI Intern -- Beaufort Solutions}{Jul. 2022 -- Nov. 2022}{St. John's, NL, Canada}
\begin{points}
\item Designed a data-collection pipeline and few-shot vision-language classifier for photo-book curation; led data collection and evaluation, contributing to \href{https://doi.org/10.1117/1.JEI.33.1.013028}{TEG (Journal of Electronic Imaging, 2024)}.
\end{points}

\role{AI Algorithm Intern -- Zhongkong Hanlian}{Dec. 2017 -- Aug. 2018}{Hangzhou, China}
\begin{points}
\item Built vehicle-logo recognition models and annotation tools in Python, C++, and Java; supported real-time deployment at community security gates.
\end{points}

\newpage
{\large\bfseries Wanglong Lu}\hfill{\small Research, Engineering \& Education}\par

\cvsection{Selected Research Publications}
\project{\href{https://doi.org/10.1109/TNNLS.2026.3707463}{UltraDiffEdit -- High-Resolution Generative Editing}}{First author; IEEE TNNLS 2026}
\begin{points}
\item Developed tuning-free editing with multi-patch encoding, global-local denoising, and progressive sampling; demonstrated \textbf{8K editing on a single RTX 3090}. \href{https://github.com/LonglongaaaGo/UltraDiffEdit}{Code}; \href{https://longlongaaago.github.io/videos/}{video}.
\end{points}

\project{\href{https://doi.org/10.1109/TVCG.2024.3434386}{FACEMUG -- Multimodal Local Facial Editing}}{First author; IEEE TVCG 2024}
\begin{points}
\item Designed multimodal feature fusion and self-supervised latent warping for local facial edits across \textbf{six input modalities}, including text, sketches, and exemplars, while preserving unedited regions.
\end{points}

\project{\href{https://neurips.cc/virtual/2025/loc/san-diego/poster/119606}{AdaMSS -- Parameter-Efficient Adaptation}}{Co-author; NeurIPS 2025}
\begin{points}
\item Contributed to multi-subspace fine-tuning; reported ViT-Large results exceeded LoRA's average accuracy across seven tasks using \textbf{15.4\% of LoRA's trainable parameters}. Method \textbf{integrated into Hugging Face PEFT}. \href{https://github.com/jzheng20/AdaMSS}{Code}; \href{https://github.com/huggingface/peft/tree/main/examples/adamss_finetuning}{PEFT integration}.
\end{points}

\project{\href{https://doi.org/10.1109/TCSVT.2026.3720433}{DocPure -- Unified Document Restoration}}{Co-first author; IEEE TCSVT 2026}
\begin{points}
\item Co-authored prompt-free document restoration using degradation-aware structure extraction and structure-guided wavelet modulation across multiple degradations. \href{https://github.com/LingmingSSS/DocPure}{Code}.
\end{points}

\project{\href{https://doi.org/10.26599/CVM.2025.9450408}{GRIG -- Data-Efficient Image Inpainting}}{First author; CVM 2025}
\begin{points}
\item Developed iterative residual inpainting using CNN features, Transformer reasoning, and forgery-patch adversarial training; evaluated on \textbf{10 few-shot datasets} using FID and LPIPS. \href{https://github.com/LonglongaaaGo/GRIG_few_shot_inpainting}{Code}; \href{https://arxiv.org/html/2304.12035v1}{benchmarks}.
\end{points}

\project{\href{https://doi.org/10.1016/j.patcog.2024.111312}{Visual Style Prompt Learning}}{First author; Pattern Recognition 2025}
\begin{points}
\item Developed diffusion-generated visual prompts and style-modulated aggregation for blind face restoration. \href{https://github.com/LonglongaaaGo/VSPBFR}{Code}.
\end{points}
{\small Full publication record: \href{https://scholar.google.com/citations?user=TuxCf4UAAAAJ}{Google Scholar} and \href{https://longlongaaago.github.io/publications/}{publication portfolio}.}

\cvsection{AI Engineering Prototypes \& Deployment}
\project{Local Agentic Code Editor}{Nasdaq internal prototype; 2026}
\begin{points}
\item Prototyped a local AI code editor connecting \textbf{language-model APIs with local tool execution}; an \textbf{agentic loop} selects code writing, code review, or shell-tool execution to support code generation and debugging.
\end{points}

\cvsection{Education}
\textbf{Ph.D., Computer Science} -- Memorial University of Newfoundland, Canada \hfill 2021--2025\par
\textbf{M.Sc., Computer Software and Theory} -- Wenzhou University, China \hfill 2018--2021\par
\textbf{B.Sc., Digital Media Technology} -- Communication University of Zhejiang, China \hfill 2014--2018

\cvsection{Selected Recognition \& Service}
Outstanding Research Award (Ph.D.), Memorial University, 2023; four granted Chinese invention patents.\par
Mentored \textbf{10 students} in generative vision and applied ML; supported industry research on cheque OCR and tabular-data generation.\par
{\raggedright Invited reviewer for \textbf{IEEE TPAMI, TIP, TMM, TCSVT}, Pattern Recognition, Knowledge-Based Systems, \mbox{Neurocomputing}, Applied Soft Computing, and ECCV 2026.\par}
Presented feature-selection research to \textbf{60+ attendees} (2026); invited speaker, IEEE Newfoundland and Labrador Section AGM (2026); guest lecturer on deep generative models, Memorial University (2025).
\end{document}
